\documentclass[11pt,a4paper]{article}

\usepackage[a4paper,margin=23mm]{geometry}
\usepackage{fontspec}
\IfFontExistsTF{Times New Roman}{\setmainfont{Times New Roman}}{\setmainfont{Arial}}
\usepackage{polyglossia}
\setmainlanguage{russian}
\newfontfamily\cyrillicfont{Times New Roman}
\newfontfamily\cyrillicfonttt{Arial}
\usepackage{amsmath,amssymb,mathtools,bm}
\usepackage{booktabs,tabularx,array}
\usepackage{xcolor}
\usepackage[colorlinks=true,linkcolor=blue!45!black,urlcolor=blue!45!black]{hyperref}
\setlength{\parindent}{0pt}
\setlength{\parskip}{0.48em}
\emergencystretch=2em

\newcommand{\R}{\mathbb{R}}
\newcommand{\A}{\mathcal{A}}
\newcommand{\ip}[2]{\left\langle #1,#2\right\rangle}
\newcommand{\norm}[1]{\left\lVert #1\right\rVert}
\newcommand{\argmin}{\operatorname*{arg\,min}}
\newcommand{\argmax}{\operatorname*{arg\,max}}
\newcommand{\dom}{\operatorname{dom}}
\newcommand{\soft}{\operatorname{soft}}

\title{\bfseries От primal ERM к dual Frank--Wolfe\\
\large Явный вывод обновлений для dual variables и исходного параметра \(w\)}
\author{Техническая записка}
\date{16 июля 2026}

\begin{document}
\maketitle

\section{Исходная задача}

Начнем с
\begin{equation}
  \label{eq:primal}
  \boxed{
  \min_{w\in\R^d}P(w)
  :=R(w)+\frac1n\sum_{i=1}^n f_i(K_iw).}
\end{equation}
Здесь \(K_i\) может быть identity; он оставлен, чтобы одна формула покрывала linear prediction, последнюю голову сети и structured score. Предполагаем, что \(R\) и \(f_i\) --- proper closed convex functions и выполнено strong duality. Для произвольной nonconvex end-to-end сети это не автоматический факт: дальнейшие формулы точны, если loss имеет convex или точное saddle/max представление.

По определению Fenchel conjugate
\[
  f_i(z)=\max_{\alpha_i}
  \left\{\ip{\alpha_i}{z}-f_i^*(\alpha_i)\right\}.
\]
Подставляем это в \eqref{eq:primal}:
\begin{align}
  P^\star
  &=\min_w\max_{\alpha_1,\ldots,\alpha_n}
  \left[
  R(w)+\frac1n\sum_i
  \bigl(\ip{\alpha_i}{K_iw}-f_i^*(\alpha_i)\bigr)
  \right].
  \label{eq:saddle}
\end{align}
Введем агрегированную dual-переменную в пространстве параметров:
\begin{equation}
  \label{eq:v}
  v(\alpha):=-\frac1n\sum_{i=1}^nK_i^*\alpha_i.
\end{equation}
После перестановки min и max минимизация по \(w\) равна
\[
  \min_w\{R(w)-\ip{v(\alpha)}{w}\}=-R^*(v(\alpha)).
\]
Поэтому dual problem имеет вид
\begin{equation}
  \label{eq:dual-max}
  \boxed{
  \max_{\alpha_i\in\dom f_i^*}
  D(\alpha):=-R^*(v(\alpha))-\frac1n\sum_i f_i^*(\alpha_i).}
\end{equation}
Для Frank--Wolfe удобнее минимизировать минус dual objective:
\begin{equation}
  \label{eq:dual-min}
  \boxed{
  \min_{\alpha\in\A}F(\alpha)
  :=R^*(v(\alpha))+\frac1n\sum_i f_i^*(\alpha_i),
  \qquad
  \A=\prod_{i=1}^n\dom f_i^*.}
\end{equation}

\section{Как из dual variables получается исходный \(w\)}

Минимизатор лагранжиана по \(w\) удовлетворяет
\[
  0\in\partial R(w)-v(\alpha).
\]
Если \(R\) строго выпукла, recovery map однозначна:
\begin{equation}
  \label{eq:recovery}
  \boxed{w(\alpha)=\nabla R^*(v(\alpha)).}
\end{equation}
Это центральная формула: мы оптимизируем \(\alpha\), а параметр исходной модели каждый раз получаем через gradient conjugate регуляризатора.

Если \(f_i^*\) дифференцируема, block-gradient dual objective равен
\begin{equation}
  \label{eq:dual-gradient}
  \boxed{
  \nabla_iF(\alpha)
  =\frac1n\left[
  \nabla f_i^*(\alpha_i)-K_iw(\alpha)
  \right].}
\end{equation}
Действительно, \(\partial v/\partial\alpha_i=-K_i^*/n\), а \(\nabla R^*(v)=w\).

\section{Frank--Wolfe update в dual-пространстве}

Пусть \(\A\) компактна и \(F\) гладкая. Полный linear minimization oracle:
\begin{equation}
  \label{eq:lmo-general}
  \boxed{
  s^t\in\argmin_{s\in\A}
  \ip{s}{\nabla F(\alpha^t)}.}
\end{equation}
Так как \(\A=\A_1\times\cdots\times\A_n\), он распадается:
\begin{equation}
  \label{eq:lmo-blocks}
  s_i^t\in\argmin_{s_i\in\A_i}
  \ip{s_i}{\nabla f_i^*(\alpha_i^t)-K_iw^t},
  \qquad i=1,\ldots,n,
\end{equation}
где общий множитель \(1/n\) отброшен. Dual Frank--Wolfe update:
\begin{equation}
  \label{eq:alpha-update}
  \boxed{
  \alpha^{t+1}=(1-\gamma_t)\alpha^t+\gamma_ts^t.}
\end{equation}

Теперь явно переведем этот update обратно в primal-параметр. Из линейности \(v(\alpha)\):
\begin{align}
  v^{t+1}
  &=(1-\gamma_t)v^t+\gamma_tv(s^t),
  \label{eq:v-update}\\
  \boxed{
  w^{t+1}
  =\nabla R^*\!\left((1-\gamma_t)v^t+\gamma_tv(s^t)\right).}
  \label{eq:w-general-update}
\end{align}
В общем случае это \emph{не} convex combination между \(w^t\) и \(w(s^t)\), потому что map \(\nabla R^*\) может быть нелинейной. Такое простое усреднение возникает для quadratic regularizer.

\subsection{Block-coordinate update}

Если обновляем только блоки \(B_t\subseteq\{1,\ldots,n\}\), то
\begin{align}
  \alpha_i^{t+1}
  &=(1-\gamma_{t,i})\alpha_i^t+\gamma_{t,i}s_i^t,
  &&i\in B_t,\\
  \alpha_i^{t+1}&=\alpha_i^t,
  &&i\notin B_t.
\end{align}
Агрегат и исходный параметр обновляются как
\begin{align}
  \boxed{
  v^{t+1}=v^t-\frac1n\sum_{i\in B_t}
  K_i^*(\alpha_i^{t+1}-\alpha_i^t),}
  \label{eq:v-block}\\
  \boxed{w^{t+1}=\nabla R^*(v^{t+1}).}
  \label{eq:w-block}
\end{align}
То есть для BCFW не нужно пересчитывать сумму по всему датасету: достаточно изменить вклад выбранных dual blocks и применить recovery map.

\section{Чистая max/hinge-форма}

Для Frank--Wolfe особенно удобен случай
\begin{equation}
  \label{eq:max-form}
  f_i(w)=\max_{a_i\in\A_i}
  \ip{a_i}{b_i-M_iw},
\end{equation}
где \(\A_i\) --- компактное множество альтернатив. Тогда
\begin{align}
  q(a)&=\frac1n\sum_iM_i^*a_i,\\
  F(a)&=R^*(q(a))-\frac1n\sum_i\ip{a_i}{b_i},
  \label{eq:max-dual}\\
  \boxed{w(a)=\nabla R^*(q(a)),}
  \label{eq:max-recovery}\\
  \nabla_iF(a)&=\frac1n(M_iw(a)-b_i).
\end{align}
Здесь LMO совсем простой:
\begin{equation}
  \label{eq:max-lmo}
  \boxed{
  s_i^t\in\argmin_{s_i\in\A_i}
  \ip{s_i}{M_iw^t-b_i}
  =\argmax_{s_i\in\A_i}
  \ip{s_i}{b_i-M_iw^t}.}
\end{equation}
Это и есть выбор наиболее нарушенной alternative/label/constraint при текущей модели. После него
\begin{align}
  a_i^{t+1}&=(1-\gamma_t)a_i^t+\gamma_ts_i^t,\\
  q^{t+1}&=q^t+\frac{\gamma_t}{n}
  \sum_{i\in B_t}M_i^*(s_i^t-a_i^t),
  \label{eq:q-update}\\
  \boxed{w^{t+1}=\nabla R^*(q^{t+1}).}
\end{align}

\section{Подстановка конкретных регуляризаторов}

Во всех формулах ниже \(q\) означает агрегат из max-form. Для общей Fenchel-формы нужно просто заменить \(q\) на \(v\).

\subsection{Обычный quadratic regularizer}

Пусть
\[
  R(w)=\frac\lambda2\norm{w}_2^2,\qquad
  R^*(q)=\frac1{2\lambda}\norm{q}_2^2.
\]
Тогда
\begin{equation}
  \boxed{w=\frac q\lambda.}
\end{equation}
Для full FW:
\begin{equation}
  \boxed{
  w^{t+1}=(1-\gamma_t)w^t
  +\gamma_t\frac{q(s^t)}{\lambda}.}
\end{equation}
Для block update:
\begin{equation}
  \boxed{
  w^{t+1}=w^t+\frac1{\lambda n}
  \sum_{i\in B_t}\gamma_{t,i}M_i^*(s_i^t-a_i^t).}
\end{equation}

\subsection{Weighted quadratic regularizer}

Для \(H\succ0\):
\[
  R(w)=\frac\lambda2w^\top Hw,\qquad
  R^*(q)=\frac1{2\lambda}q^\top H^{-1}q,\qquad
  \boxed{w=\frac1\lambda H^{-1}q.}
\]
Следовательно,
\begin{equation}
  \boxed{
  w^{t+1}=w^t+\frac1{\lambda n}H^{-1}
  \sum_{i\in B_t}\gamma_{t,i}M_i^*(s_i^t-a_i^t).}
\end{equation}

\subsection{\(p\)-power regularizer}

Пусть \(p>1\), \(p^{-1}+r^{-1}=1\), и
\[
  R(w)=\frac\lambda p\norm{w}_p^p.
\]
Тогда
\[
  R^*(q)=\frac{\lambda^{1-r}}r\norm{q}_r^r,\qquad
  \boxed{
  w_j=\lambda^{1-r}\operatorname{sign}(q_j)|q_j|^{r-1}.}
\]
После FW update сначала обновляется \(q^{t+1}\) по \eqref{eq:q-update}, затем эта формула применяется покоординатно. При \(p\ne2\) update по \(w\) нелинеен.

\subsection{Elastic net}

Пусть
\[
  R(w)=\lambda_1\norm{w}_1+\frac{\lambda_2}{2}\norm{w}_2^2,
  \qquad \lambda_2>0.
\]
Тогда
\begin{align}
  R^*(q)&=\frac1{2\lambda_2}
  \norm{\soft(q,\lambda_1)}_2^2,\\
  \boxed{
  w=\frac1{\lambda_2}\soft(q,\lambda_1),}
\end{align}
где \(\soft(q,\tau)_j=\operatorname{sign}(q_j)(|q_j|-\tau)_+\). Поэтому BCFW обновляет \(q\) линейно, а затем применяет soft-thresholding.

\subsection{Entropy regularizer на simplex}

Если \(w\in\Delta_d\) и
\[
  R(w)=\lambda\sum_jw_j\log w_j+I_{\Delta_d}(w),
\]
то
\begin{align}
  R^*(q)&=\lambda\log\sum_j\exp(q_j/\lambda),\\
  \boxed{w=\operatorname{softmax}(q/\lambda).}
\end{align}
Здесь FW update dual aggregate \(q\) превращается в multiplicative/softmax-like изменение исходного \(w\), а не в линейное усреднение.

\subsection{Чистый \(\ell_1\)-regularizer}

Для \(R(w)=\lambda\norm{w}_1\):
\[
  R^*(q)=I_{\{\norm{q}_\infty\le\lambda\}}(q).
\]
Conjugate негладкая, а \(w=\nabla R^*(q)\) не определено однозначно. Поэтому чистый \(\ell_1\) не дает описанный выше гладкий dual FW без дополнительного smoothing или небольшой quadratic component. Elastic net устраняет эту проблему.

\section{Итоговая схема}

Для max/hinge DL subproblem вся итерация сводится к четырем строкам:
\begin{align*}
  w^t&=\nabla R^*(q^t),\\
  s_i^t&\in\argmax_{s_i\in\A_i}\ip{s_i}{b_i-M_iw^t},
  \qquad i\in B_t,\\
  q^{t+1}&=q^t+\frac1n\sum_{i\in B_t}
  \gamma_{t,i}M_i^*(s_i^t-a_i^t),\\
  a_i^{t+1}&=(1-\gamma_{t,i})a_i^t+\gamma_{t,i}s_i^t,
  \qquad
  w^{t+1}=\nabla R^*(q^{t+1}).
\end{align*}
Регуляризатор полностью определяет отображение \(q\mapsto w\). Quadratic дает линейный update, \(p\)-regularizer --- степенной, elastic net --- soft-thresholding, entropy --- softmax.

\end{document}
